Federated Learning (FL) enables collaborative model training across distributed data sources without sharing raw data, preserving user privacy~\citep{fed-learn, fed-learn-adv, fed-learn-survey}.
However, data heterogeneity among users remains a critical challenge, leading to suboptimal global model performance for individual users.
Various approaches have been proposed to address this, ranging from fine-tuning and meta-learning adaptations to fully personalized federated learning (PFL) techniques.

A simple yet effective strategy for personalization is fine-tuning, where the global model is adjusted using local user data post-training~\citep{ft-is-fine, pfl-bench, pfl}.
It is necessary to clarify that we consider fine-tuning separate from PFL solutions, as the additional computation and models happen after the federated training is completed.
While computationally efficient during federated training, fine-tuning assumes access to sufficient local data, and its success largely depends on the quality of the shared initialization.
Multiple works have found that FedAvg with fine-tuning typically outperforms more recent meta-learning and PFL methods~\citep{ft-is-fine, benchmark-for-pfl, bridging-pfl}.

Meta-learning methods extend this idea by learning a globally shared model well-suited for rapid personalization~\citep{maml, meta-learning1, meta-learning2, debias-ml}.
In the context of FL, FOMAML~\citep{maml} and Per-FedAvg~\citep{per-fedavg} have been adapted to optimize a shared initialization that accelerates convergence during fine-tuning.
In addition to applying conventional meta-learning algorithms to the federated learning context, a variant named Reptile \citep{reptile} can be shown to be equivalent to FedAvg under the condition of equally sized local data sets.
In this article, we exclusively consider aggregation with equally weighted users.
Therefore, we will use the term FedAvg to encompass both FedAvg and Reptile and refer to FedAvg as a meta-learning method.
Meta-learning approaches often have limitations, including the requirement for additional computation for second-order gradients and the need for fine-tuning for the initialization to perform well.
More recent adaptations of meta-learning for FL include FedABML~\citep{fedabml}.
However, in their experiments, FedABML typically performs around a single percentage point better than FedAvg with fine-tuning.
% pFedMe compared against Per-FedAvg

Next, PFL approaches address data diversity by learning personalized models for each user during federated training~\citep{toward-pfl, pfl}.
Other personalization techniques than meta-learning include clustering~\citep{cluster-plus, cluster-sattler, cluster-briggs}, model mixture~\citep{ditto, cluster-plus, fedem, pfl-fo}, parameter decoupling~\citep{fedbn, pfl-layers, shared-pfl, think-locally}, and knowledge distrillation~\citep{ensemble-distil, fedme, dfml, mutual-fl}.
These intricate solutions often include learning a distinct model for each user throughout training \citep{ditto}, sharing only a subset of the entire model globally \citep{fedbn}, or treating user data as a mixture of distributions \citep{fedem}.
In general, PFL methods often require substantial increases in memory and computation to produce and communicate additional models during federated training, which challenges scalability and resource efficiency.
Furthermore, several prevailing methods may prove restrictive or inequitable for applications accommodating new users.
For instance, Ditto's approach \citep{ditto} of training personalized models during federated learning becomes infeasible for users who do not participate in training.
Likewise, FedBN \citep{fedbn} and FedEM \citep{fedem}, which rely on insights from user data, encounter difficulties when new users lack sufficient data for tailoring a model to new users.

Importantly, this work focuses on how including within-round learning rate decay can improve FL performance without significant extra computation or the need for additional models.
Moreover, unlike the aforementioned meta-learning applications to FL, the learning rate decay can flexibly emphasize initial model performance and rapid personalization.
We demonstrate that within-round learning rate decay is a generalization of FedAvg or Reptile, making them particularly relevant baselines for comparison.
Several other works propose to improve the original FedAvg procedure.
SCAFFOLD~\citep{scaffold} reduces client drift using control variates; FedProx~\citep{fedprox} adds a proximal penalty to the local objective to limit divergence from the global model; and FedAdam and FedYogi~\citep{fed-opt} apply adaptive moment estimation at the server during aggregation.
Because these methods modify the global aggregation step or augment the local objective rather than the within-round learning rate schedule, they are orthogonal to FedDecay and could in principle be combined with our approach for further gains.
Unlike our work, they do not exploit the local learning rate within a communication round to control the balance between initial model success and rapid personalization.
While we do not propose a PFL method by definition, since fine-tuning occurs after federated training, we include PFL methods in our comparisons to assess the performance gap between our process and PFL techniques.
More recent PFL solutions exist, like \citet{knn-per}, but comparison with PFL solutions is not the main focus of our work.
More importantly, experiments demonstrate consistent improvements compared to FedAvg and FOMAML and close the performance gap with PFL methods with a simple, interpretable learning rate modification without additional computational or memory overheads.
